\documentclass[10pt,a4paper]{amsart}
\usepackage[margin=19mm]{geometry}
\usepackage{amsmath,amssymb,mathtools}
\usepackage[scr=rsfs,cal=euler]{mathalfa}
\usepackage{xfrac}
\usepackage[unicode,hidelinks,bookmarks=false]{hyperref}
\newcommand{\ie}{i.e.\ }
\newcommand{\cf}{cf.\ }
\newcommand{\resp}{respectively\ }
\newcommand{\Subsection}{\subsection}
\providecommand{\nobreakdash}{\nobreak\hbox{-}}
\setlength{\emergencystretch}{2em}
\multlinegap0pt
\usepackage{mathtools}
\usepackage{amssymb}
\usepackage{braket}
\usepackage{dsfont}
\usepackage{enumerate}
\usepackage{tikz}
\usepackage{algorithm}\usepackage{algpseudocode}
\newtheorem{lemma}{Lemma}
\newtheorem{theorem}{Theorem}
\usepackage{cleveref}
\crefname{lemma}{Lemma}{Lemmas}
\crefname{theorem}{Theorem}{Theorems}
\newcommand{\IOprod}[1]{\left( #1 \right)_{I\times\Omega}}
\newcommand{\ImOprod}[1]{\left( #1 \right)_{I_m\times\Omega}}
\newcommand{\LtwoHone}{L^2(I;H^1(\Omega))}
\newcommand{\N}{\mathbb{N}}
\newcommand{\Oprod}[1]{\left( #1 \right)_{\Omega}}
\newcommand{\R}{\mathbb{R}}
\newcommand{\Ukh}{U_{kh}}
\newcommand{\Vkh}{V_{kh}}
\newcommand\blue[1]{\textcolor{black}{#1}}
\newcommand{\cO}[1]{c\left( #1 \right)}
\newcommand{\chatIO}[1]{\hat c \left(\!\left( #1 \right)\!\right)}
\newcommand{\chatO}[1]{\hat c\left( #1 \right)}
\newcommand{\ekh}{e_{kh}}
\newcommand{\pkh}{p_{kh}}
\newcommand{\qkh}{q_{kh}}
\newcommand{\ukh}{u_{kh}}
\newcommand{\vk}{v_{k}}
\newcommand{\vkh}{v_{kh}}
\title{Error estimates for finite element discretizations of the instationary Navier-Stokes equations}
\author{Boris Vexler, Jakob Wagner}
\date{Source: arXiv:2307.14217v2 (CC BY 4.0)}
\begin{document}
\maketitle

\noindent\emph{Source and licence.} Error estimates for finite element discretizations of the instationary Navier-Stokes equations. Version arXiv:2307.14217v2 (CC BY 4.0), distributed by arXiv under the Creative Commons Attribution 4.0 International licence, http://creativecommons.org/licenses/by/4.0/, as recorded in the arXiv licence metadata for this exact version and shown on the abstract page https://arxiv.org/abs/2307.14217v2. Full licence notice: \texttt{licenses/arxiv-2307.14217-CC-BY-4.0.txt}.\par\smallskip

\noindent\textit{Editorial scope.} Counted unit: the article's theorem thm:discrete\_nav\_stokes\_unique (source0.tex lines 1960--1963). The article's complete original proof of it is delivered with it (source0.tex lines 1964--2033, one \texttt{proof} environment of 70 lines). No mathematics is written, repaired or re-derived: the statement and the proof are the article's own lines, in the article's own notation, and no retained line is substituted or re-flowed.  Retained uncounted as the source-local context the proof needs: lines 488--509; lines 918--924; lines 1097--1105; lines 1280--1288; lines 1300--1303; lines 1310--1312; lines 1433--1446; lines 1557--1566; lines 1778--1788; lines 1930--1938.  Nothing was omitted from any retained span, so the package carries no omission record.  No retained line contains a citation, so the wrapper carries no bibliography and no bibliography entry.  No retained line contains a figure environment, an image-inclusion command or drawing code, so no image is delivered and the package declares no asset dependency.  Editorial formatting only, in addition: an AMS \texttt{amsart} preamble that restates the article's own declarations and macros used by the retained lines, namely the environments \texttt{lemma}, \texttt{theorem} on separate counters, together with the standard packages those lines need.  The retained environments print in this document as Lemma 1, Theorem 1 in order of appearance, whereas the article prints the same items with the numbering of its own full document.  The complete native source of the article as distributed in the arXiv e-print for this exact version is bundled unchanged as \texttt{sources/144815/source0.tex}.
\begin{lemma}\label{lemm:c_swap}
  Let $\Omega \subset \R^2$ be an open Lipschitz domain, then there holds the estimate
  \begin{align}
    % \|v\|_{L^4(\Omega)} &\le 2^{-\frac{1}{4}}\|v\|_{L^2(\Omega)}^{\frac{1}{2}}\|\nabla v\|_{L^2(\Omega)}^{\frac{1}{2}}
    % && \text{for all } v \in H^1_0(\Omega),
    % \label{eq:l4_interpol_h10}
    % \\
    \|v\|_{L^4(\Omega)} &\le C\|v\|_{L^2(\Omega)}^{\frac{1}{2}}\|v\|_{H^1(\Omega)}^{\frac{1}{2}}
    \qquad \text{for all } v \in H^1(\Omega),
    \label{eq:l4_interpol_general}
  \end{align}
  due to which, the trilinear form $\cO{\cdot,\cdot,\cdot}$ satisfies for all $u,v,w \in H^1_0(\Omega)^2$:
  \begin{align*}
    \cO{u,v,w} & \le \|u\|_{L^4(\Omega)} \|\nabla v\|_{L^2(\Omega)} \|w\|_{L^4(\Omega)},\\
    \cO{u,v,u} & \le C\|u\|_{L^2(\Omega)} \|\nabla u\|_{L^2(\Omega)}\|\nabla v\|_{L^2(\Omega)}.
  \end{align*}
  Let further $\nabla \cdot u = 0$. Then it holds
  \begin{equation*}
    \cO{u,v,w} = - \cO{u,w,v} \quad \text{and} \quad 
    \cO{u,v,v} = 0.
  \end{equation*}
\end{lemma}

  \begin{equation}\label{eq:nav_stokes_weak_with_pressure}
    \blue{\IOprod{\partial_t u,v} }
    + \nu \IOprod{\nabla u, \nabla v} 
    + \IOprod{(u\cdot \nabla)u,v}
    - \IOprod{p,\nabla \cdot v} + \IOprod{q, \nabla \cdot u} 
    = \IOprod{f,v},
  \end{equation}

\begin{lemma}\label{lemm:chat}
  The trilinear form $\chatO{\cdot,\cdot,\cdot}$ satisfies
  \begin{align*}
    \chatO{u_h,v_h,w_h} & \le C\|\nabla u_h\|_{L^2(\Omega)} \|\nabla v_h\|_{L^2(\Omega)} 
    \|\nabla w_h\|_{L^2(\Omega)} && \text{for all } u_h, v_h, w_h \in U_h,\\
    \chatO{u_h, v_h, w_h} &= - \chatO{u_h,w_h,v_h} && \text{for all } u_h, v_h, w_h \in U_h,\\
    \chatO{u_h, v_h, v_h} &= 0 && \text{for all } u_h, v_h \in U_h.
  \end{align*}
\end{lemma}

\begin{lemma}\label{lemm:jump_reorder}
  For any $\vk\in X_k^q(L^2(\Omega))$ it holds
  \begin{align*}
  \ImOprod{\partial_t \vk,\vk} + \Oprod{[\vk]_{m-1},v_{kh,m-1}^+}
  &= \frac{1}{2} \left(\|v_{k,m}^-\|_{L^2(\Omega)}^2  + \|[\vk]_{m-1}\|_{L^2(\Omega)}^2 - \|v_{k,m-1}^-\|_{L^2(\Omega)}^2 \right),\\
  - \ImOprod{\vk, \partial_t \vk} - \Oprod{v_{k,m}^-, [\vk]_{m}}
  &= \frac{1}{2} \left(\|v_{k,m-1}^+\|_{L^2(\Omega)}^2  + \|[\vk]_{m}\|_{L^2(\Omega)}^2 - \|v_{k,m}^+\|_{L^2(\Omega)}^2 \right).
  \end{align*}
\end{lemma}

\begin{equation}\label{eq:nav_stokes_discrete}
  \mathfrak B (\ukh,\vkh) + \chatIO{\ukh,\ukh,\vkh} = \IOprod{f,\vkh} + \Oprod{u_0,v_{kh,0}^+} 
  \quad \text{for all } \vkh \in \Vkh.
\end{equation}

\begin{equation}\label{eq:nav_stokes_discrete_with_pressure}
  B((\ukh,\pkh),(\vkh,\qkh)) + \chatIO{\ukh,\ukh,\vkh} = \IOprod{f,\vkh} + \Oprod{u_0,v_{kh,0}^+},
\end{equation}

\begin{lemma}\label{thm:discrete_gronwall_quadlinear}
  Let $\{x_n\}, \{b_n\}, \{c_n\}, \{d_n\}, \{\gamma_n\}$ be sequences of nonnegative numbers
  and $B \ge 0$ a constant, such that for each $n \in \N_0$ it holds
  \begin{equation*}
    x_n^2 + \sum_{m=0}^n b_m \le \sum_{m=0}^n \gamma_m x_m^2 
    + \sum_{m=0}^n d_m x_m
    + \sum_{m=0}^n c_m + B,
  \end{equation*}
  and $\gamma_m < 1$ for all $m \in \{0,...,n\}$. Then with $\sigma_m := (1 - \gamma_m)^{-1}$ it holds
  \begin{equation*}
    x_n^2 + \sum_{m=0}^n b_m \le 2\exp \left(2 \sum_{m=0}^n \sigma_m \gamma_m \right)\cdot
    \left[\left(\sum_{m=0}^n d_m \right)^2 + \sum_{m=0}^n c_m + B\right].
  \end{equation*}
\end{lemma}

\begin{theorem}\label{thm:stability_discrete_navier_stokes}
  Let $f \in L^1(I;L^2(\Omega)^2) + L^2(I;\blue{H^{-1}(\Omega)^2})$, $u_0 \in H$ and $u_{kh} \in \Vkh$ satisfy
  \cref{eq:nav_stokes_discrete} for either $q=0$ or $q=1$. Then there holds the bound
  \begin{equation*}
    \|\ukh\|_{L^\infty(I;L^2(\Omega))} + \blue{\sqrt{\nu}} \|\ukh\|_{\blue{\LtwoHone}} 
    \le C \left(\|u_0\|_H + \|f\|_{L^1(I;L^2(\Omega)) + L^2(I;\blue{H^{-1}(\Omega)})}\right),
  \end{equation*}
  \blue{with a constant $C$ only depending on $\Omega,\nu$. If $f \in L^1(I;L^2(\Omega)^2)$, then 
  $C$ only depends on $\Omega$.}
\end{theorem}

\begin{theorem}\label{thm:nav_stokes_l2h1}
  Let $f \in L^2(I;L^2(\Omega)^2)$, $u_0 \in V$ and let $(u,p)$, $(\ukh,\pkh)$ \blue{solve the} continuous and
  fully discrete Navier-Stokes equations \eqref{eq:nav_stokes_weak_with_pressure} and
\eqref{eq:nav_stokes_discrete_with_pressure} respectively.
\blue{Further let $k$ be small enough.}
Then for any $\chi_{kh} \in \Vkh$, there holds 
\begin{align*}
  \|\nabla(u-\ukh)\|_{L^2(I\times\Omega)} & \le C \left(\| \nabla(u - \chi_{kh})\|_{L^2(I\times\Omega)} 
   + \|\nabla(u - \pi_\tau u)\|_{L^2(I\times \Omega)} + \|\nabla(u - R_h^S(u,p))\|_{L^2(I\times \Omega)}\right).
\end{align*}
\end{theorem}

\begin{lemma}\label{lemm:vanishing_support_u}
  Let $u \in L^2(I;H^1_0(\Omega)^2)$, and $u_{kh} \in \Ukh $ such that
  $\|u - u_{kh}\|_{L^2(I;H^1(\Omega))} \to 0$ as $(k,h) \to 0$. Then it holds
  \begin{align*}
    \sup_{1 \le m \le M} \int_{I_{m,k}} \|\nabla u\|_{L^2(\Omega)}^2 \ dt  \xrightarrow{(k,h) \to 0} 0
    \quad \text{and} \quad
    \sup_{1 \le m \le M} \int_{I_{m,k}} \|\nabla u_{kh}\|_{L^2(\Omega)}^2 \ dt \xrightarrow{(k,h) \to 0} 0.
  \end{align*}
\end{lemma}

  \begin{theorem}\label{thm:discrete_nav_stokes_unique}
    Let $f \in L^2(I;L^2(\Omega)^2)$ and $u_0 \in V$. Then for $(k,h)$ small enough, the solution
    $\ukh$ to \eqref{eq:nav_stokes_discrete} is unique.
  \end{theorem}

  \begin{proof}
    \blue{To show uniqueness, we assume two solutions $\ukh^1, \ukh^2$ and define $\ekh:= \ukh^1 - \ukh^2$.
  Testing the equations for $m=1,...,M$ with $\ekh \chi_m$, where $\chi_m$ denotes the characteristic
  function of the interval $I_m$, and subtracting them leads to
  \begin{equation*}
    \mathfrak B (\ekh,\ekh \chi_m) + \chatIO{\ukh^1,\ukh^1,\ekh \chi_m} - \chatIO{\ukh^2,\ukh^2,\ekh \chi_m} = 0
  \end{equation*}
  Adding and subtracting $\chatIO{\ukh^1,\ukh^2,\ekh \chi_m}$ and applying \Cref{lemm:chat} yields
  \begin{equation*}
    \mathfrak B (\ekh,\ekh \chi_m) + \chatIO{\ekh,\ukh^2,\ekh \chi_m}= 0.
  \end{equation*}
  Applying \Cref{lemm:jump_reorder}, the bilinear form can be written  as 
  \begin{equation*}
    \mathfrak B (\ekh,\ekh \chi_m) = 
    \frac{1}{2} \left(\|e_{kh,m}^-\|_{L^2(\Omega)}^2  
    + \|[\ekh]_{m-1}\|_{L^2(\Omega)}^2 - \|e_{kh,m-1}^-\|_{L^2(\Omega)}^2 \right)
    + \nu\|\nabla \ekh\|_{L^2(I_m;L^2(\Omega))}^2.
  \end{equation*}
  From the definition of $\hat c$, the estimates of \Cref{lemm:c_swap}, and Hölder's inequality in time,
  we obtain
  \begin{align*}
    \chatIO{\ekh,\ukh^2,\ekh \chi_m} &=
    \frac{1}{2} \left[\ImOprod{(\ekh \cdot \nabla)\ukh^2,\ekh} 
        - \ImOprod{(\ekh \cdot \nabla)\ekh,\ukh^2}\right]\\
                                     &\le C \|\ekh\|_{L^\infty(I_m;L^2(\Omega))}\|\nabla \ekh\|_{L^2(I_m;L^2(\Omega))}\| \nabla \ukh^2\|_{L^2(I_m;L^2(\Omega))}\\
    & \quad + C \|\ekh\|_{L^\infty(I_m;L^2(\Omega))}^{\frac{1}{2}}
      \|\nabla \ekh\|_{L^2(I_m;L^2(\Omega))}^{\frac{3}{2}}
        \|\ukh^2\|_{L^\infty(I_m;L^2(\Omega))}^{\frac{1}{2}} 
        \| \nabla \ukh^2\|_{L^2(I_m;L^2(\Omega))}^{\frac{1}{2}}.
  \end{align*}
  We can apply Young's inequality and obtain
  \begin{align*}
    \chatIO{\ekh,\ukh^2,\ekh \chi_m} &\le
    \frac{\nu}{2}\|\nabla \ekh\|_{L^2(I_m;L^2(\Omega))}^2
    + C \|\ekh\|_{L^\infty(I_m;L^2(\Omega))}^2 \left(1 + \|\ukh^2\|_{L^\infty(I_m;L^2(\Omega))}^2\right)
    \| \nabla \ukh^2\|_{L^2(I_m;L^2(\Omega))}^2.
  \end{align*}
  where the first summand can be absorbed into $\mathfrak B(\ekh,\ekh \chi_m)$.
  Thus on each time interval, there holds the bound
  \begin{align*}
    \frac{1}{2} \left(\|e_{kh,m}^-\|_{L^2(\Omega)}^2  
    + \|[\ekh]_{m-1}\|_{L^2(\Omega)}^2 - \|e_{kh,m-1}^-\|_{L^2(\Omega)}^2 
    + \nu\|\nabla \ekh\|_{L^2(I_m;L^2(\Omega))}^2\right)\\
    \le C \|\ekh\|_{L^\infty(I_m;L^2(\Omega))}^2 \left(1 + \|\ukh^2\|_{L^\infty(I_m;L^2(\Omega))}^2\right)
    \| \nabla \ukh^2\|_{L^2(I_m;L^2(\Omega))}^2.
  \end{align*}
  Multiplying the derived inequality by 2 and summing up the inequalities for $m=1,...,n$ yields 
  for any $n \in \{1,...,M\}$
  \begin{align*}
    \|e_{kh,n}^-\|_{L^2(\Omega)}^2 + \sum_{m=1}^n 
    \left(\|[\ekh]_{m-1}\|_{L^2(\Omega)}^2 + \nu \|\nabla \ekh\|_{L^2(I_m;L^2(\Omega))}^2\right)\\
    \le 
    C \sum_{m=1}^n \|\ekh\|_{L^\infty(I_m;L^2(\Omega))}^2 
    \left(1 + \|\ukh^2\|_{L^\infty(I_m;L^2(\Omega))}^2\right) 
    \| \nabla \ukh^2\|_{L^2(I_m;L^2(\Omega))}^2.
  \end{align*}
  Here we have used $e_{kh,0}^- = 0$ as $u_{kh}^1$ and $u_{kh}^2$ satisfy the same initial condition.
  As in the proof of \Cref{thm:stability_discrete_navier_stokes}, 
  $\|\ekh\|_{L^\infty(I_m;L^2(\Omega))}$ can be bounded by evaluations of one-sided limits at the time nodes.
  Due to $f \in L^2(I;L^2(\Omega)^2)$ and $u_0 \in V$, \Cref{thm:nav_stokes_l2h1} holds,
  and thus \Cref{lemm:vanishing_support_u} yields $\| \nabla \ukh^2\|_{L^2(I_m;L^2(\Omega))}^2 \to 0$
  uniformly in $m$ for $(k,h)\to 0$. Together with \Cref{thm:stability_discrete_navier_stokes}, bounding 
  $\|\ukh^2\|_{L^\infty(I_m;L^2(\Omega))}$, we can choose the discretization fine enough, such that
  $C (1 + \|\ukh^2\|_{L^\infty(I_m;L^2(\Omega))}^2)\| \nabla \ukh^2\|_{L^2(I_m;L^2(\Omega))}^2< \frac{1}{2}$
  for all $m$.
  All in all, as in \Cref{thm:stability_discrete_navier_stokes}, an application of 
  the Gronwall \Cref{thm:discrete_gronwall_quadlinear} shows that $\ekh = 0$,
  concluding the proof of uniqueness.
  }
  \end{proof}

\end{document}
